\chapter{讨论、结论与后续建议}
\label{chap:discussion}

\section{当前实现的优点}
从工程角度看，这个项目最突出的优点是职责边界清楚。主脚本把配置、流程控制和后处理放在 MATLAB 层，保证了阅读体验；高频粒子对循环则被压到 MEX 中，保证了性能。当前版本进一步把 MATLAB 侧整理为单主文件壳层，并把长段推进与监控逻辑收敛到 \code{advance\_shell\_step} 和 \code{wall\_shear\_monitor} 两个高层接口，因此既保留了 MATLAB 层的可读性，也减少了主循环里直接拼接底层算子的负担。

第二个优点是边界处理路径已经显著简化。最近一次提交删除了 \code{bc\_mode}、PI 控制和 Neumann/ghost 分支，把旧的多套壁面逻辑收敛到单层 shell 无滑移实现。这不仅减少了代码复杂度，也明显降低了“文档写一套、代码跑另一套”的风险。

第三个优点是验证链条仍然完整。主程序能输出稳态剖面和速度场，中截面演化图能展示从静止到稳态的建立过程，运行期切应力日志则补充了近壁物理量监测。对学术写作来说，这意味着仓库不仅能支撑“我实现了什么”，还能支撑“我为什么相信它是对的”。

\section{当前实现的局限}
局限也同样明显。第一，当前报告虽然已经与源码同步到本轮重构后的状态，但文档内部仍然包含大量固定行号和 mode 列表。随着项目继续演化，这类信息比公式本身更容易过时，若不及时更新，会直接影响后续论文写作和对外交流。

第二，当前统一时间步虽然实现简单，但在更高 Reynolds 数或更复杂驱动下可能偏保守。对于当前基准算例这不是问题；但若后续扩展到长时间非稳态问题，就需要重新评估单层 CFL 是否仍然是最合适的时间推进策略。

第三，\code{advance\_shell\_step} 当前是“高层门面复用底层 mode”的实现，结构更清晰了，但内部仍存在一定的 \code{mxArray} 编排开销；如果后续把它进一步下沉为共享 core，性能和可维护性还有继续提升的空间。第四，邻居搜索 MEX 仍未真正并行化，而物理 MEX 已经在多个核心循环中使用了 OpenMP。随着粒子规模继续提升，性能瓶颈有可能逐渐转移到邻居搜索上。第五，壁面切应力目前只打印到日志，没有结构化写入文件，因此后续做批量对比时还不够方便。

\section{对文档和代码同步的建议}
如果只允许做最小改动，我更建议优先同步文档口径与输出接口，而不是立刻再做新一轮算法重构。具体来说：
\begin{enumerate}[leftmargin=2.5em]
    \item 在报告与附录中持续同步主脚本行号、MEX mode 列表和代码索引；
    \item 若继续追求性能，可把 \code{advance\_shell\_step} 从门面复用模式继续下沉为共享核心实现；
    \item 为壁面切应力、时间步和收敛统计增加结构化输出，便于后续做参数扫描；
    \item 若粒子规模继续增大，再优先评估邻居搜索并行化。
\end{enumerate}
这些工作对数值结果几乎没有影响，却会显著降低后续维护和学术写作的沟通成本。

\section{结论}
本报告对 \code{SPH-Poiseuille-Flow} 仓库做了一次从物理到实现的系统梳理。连续层面上，它求解的是二维平板泊肃叶流，驱动力由目标平均速度反推；离散层面上，它采用 cubic spline 核、cell-linked list 邻居搜索、密度重初始化、核梯度修正和弱可压状态方程；算法层面上，它采用统一 CFL 时间步与两阶段积分；边界层面上，它采用周期边界、单层 shell 壁面、无滑移接触和防穿透；实现层面上，它把高层调度保留在单个 MATLAB 壳层文件中，并通过 \code{advance\_shell\_step} 与 \code{wall\_shear\_monitor} 两个高层 MEX 接口承接长段计算。

现有结果表明，该仓库已经能够稳定再现泊肃叶流的解析速度剖面，并用速度演化与壁面切应力日志补充了运行期验证。换句话说，这个项目不只是一个“能跑出图”的脚本集合，而是一个结构清晰、版本口径统一、可继续扩展的数值基准算例基础。

\section{后续工作}
若把这个仓库作为后续研究平台，建议按以下顺序扩展：先保持报告、README 与代码索引同步，再补充更多运行期统计量；随后把壁面监测结果输出成结构化数据，便于做参数扫描和误差对比；再把 \code{advance\_shell\_step} 的内部实现从门面复用继续下沉；最后再考虑把邻居搜索并行化、把 shell 壁面接口进一步模块化，逐步向更完整的 SPH 基准平台推进。
